\documentclass[11pt,a4paper]{article}
\usepackage[utf8]{inputenc}
\usepackage[T1]{fontenc}
\usepackage[spanish,es-minimal,es-noshorthands]{babel}
\AtBeginDocument{\renewcommand{\sin}{\operatorname{sen}}\renewcommand{\arcsin}{\operatorname{arcsen}}}
\usepackage{lmodern}
\usepackage[sfdefault,lining]{FiraSans}
\usepackage[a4paper,top=2.3cm,bottom=2.2cm,left=1.8cm,right=1.8cm,headheight=1.3cm,headsep=0.4cm,footskip=1cm]{geometry}
\usepackage{amsmath,amssymb}
\usepackage{icomma}
\usepackage{graphicx}
\usepackage[table]{xcolor}
\usepackage{tikz}
\usepackage{fontawesome5}
\usepackage{ccicons}
\usepackage{fancyhdr}
\usepackage{lastpage}
\usepackage{tasks}
\usepackage{needspace}
\usepackage[most]{tcolorbox}
\usepackage{array}
\usepackage[hidelinks]{hyperref}
\AtBeginDocument{\hypersetup{pdftitle={Ficha de ejercicios de la unidad 9: Cuerpos geométricos. Áreas y volúmenes. 2º ESO},pdfsubject={Matemáticas 2º ESO},pdfkeywords={poliedros, fórmula de Euler, poliedros regulares, prisma, pirámide, cilindro, cono, esfera, área, volumen, capacidad}}}

\definecolor{navy}{HTML}{1B2748}
\definecolor{accent}{HTML}{E8603A}
\definecolor{teal}{HTML}{1F9E90}
\definecolor{soft}{HTML}{EEF1F8}
\definecolor{muted}{HTML}{5D6685}

\setlength{\parindent}{0pt}
\setlength{\parskip}{3pt}
\settasks{label=\alph*), label-format=\bfseries\color{navy}, label-width=1.3em, item-indent=1.8em, column-sep=0.6em, after-item-skip=5pt, before-skip=3pt, after-skip=4pt}

% ---------- capçalera i peu ----------
\newcommand{\llicencia}{\href{https://creativecommons.org/licenses/by-nc-sa/4.0/deed.es}{\ccbyncsa}\ \ Licencia Creative Commons Atribución-NoComercial-CompartirIgual 4.0 Internacional (CC BY-NC-SA 4.0)}
\pagestyle{fancy}
\fancyhf{}
\fancyhead[L]{\raisebox{-0.35\height}{\includegraphics[height=0.95cm]{logo-gva.png}}}
\fancyhead[C]{\small\color{muted}Ficha de ejercicios · Unidad 9: Cuerpos geométricos. Áreas y volúmenes · 2\textsuperscript{o} ESO}
\fancyhead[R]{\raisebox{-0.35\height}{\includegraphics[height=0.95cm]{logo-vermellar.png}}}
\fancyfoot[L]{\scriptsize\color{muted}\llicencia}
\fancyfoot[R]{\small\color{muted}\thepage}
\renewcommand{\headrulewidth}{0.4pt}
\renewcommand{\headrule}{\color{muted!40}\hrule height \headrulewidth}
\fancypagestyle{primera}{\fancyhf{}\renewcommand{\headrulewidth}{0pt}
  \fancyfoot[L]{\scriptsize\color{muted}\llicencia}
  \fancyfoot[R]{\small\color{muted}\thepage}}

% ---------- nivells i exercicis ----------
\newcommand{\estrelles}[1]{{\color{accent}\ifcase#1\or\faStar\faStar[regular]\faStar[regular]\or\faStar\faStar\faStar[regular]\or\faStar\faStar\faStar\fi}}
\newcounter{exer}
\newcommand{\ex}[1]{\par\needspace{4\baselineskip}\medskip\stepcounter{exer}\noindent{\color{navy}\bfseries\theexer.}\ \ {\small\estrelles{#1}}\ \ }

\newcommand{\seccio}[2]{\par\needspace{10\baselineskip}\bigskip
  \begin{tcolorbox}[colback=navy,colframe=navy,coltext=white,arc=2mm,boxrule=0pt,left=3mm,right=3mm,top=1.2mm,bottom=1.2mm]
  {\large\bfseries #1\quad #2}
  \end{tcolorbox}\vspace{-2pt}}

\newcolumntype{C}[1]{>{\centering\arraybackslash}p{#1}}
\usepackage{pifont}
% quadre de recordatori (només a la fitxa d'exercicis)
\newcommand{\recorda}[1]{\par\smallskip\begin{tcolorbox}[colback=accent!7,colframe=accent!55,arc=2mm,boxrule=0.5pt,left=2.5mm,right=2.5mm,top=1mm,bottom=1mm]\small{\color{accent}\faLightbulb[regular]}\ \textbf{Recuerda:}\ #1\end{tcolorbox}}
\newcommand{\res}[1]{#1}\newcommand{\rest}[1]{#1}

\def\fmneed{14}
% =====================================================================
% Capçalera FentMates de les fitxes: un color per a cada curs.
% Els colors són els de la pàgina inicial (LEVELS a assets/js/materials.js);
% el logo ve de _eines/logo/pdf/ (python3 _eines/logo/genera.py).
%
% Ús, al final del preàmbul de la unitat:
%   % =====================================================================
% Capçalera FentMates de les fitxes: un color per a cada curs.
% Els colors són els de la pàgina inicial (LEVELS a assets/js/materials.js);
% el logo ve de _eines/logo/pdf/ (python3 _eines/logo/genera.py).
%
% Ús, al final del preàmbul de la unitat:
%   % =====================================================================
% Capçalera FentMates de les fitxes: un color per a cada curs.
% Els colors són els de la pàgina inicial (LEVELS a assets/js/materials.js);
% el logo ve de _eines/logo/pdf/ (python3 _eines/logo/genera.py).
%
% Ús, al final del preàmbul de la unitat:
%   \input{../comu/fentmates.tex}
%   \fmcurs{1eso}
% i, al començament del document:
%   \capcalera{Matemàtiques}{1r ESO}{Fitxa d'exercicis · Unitat 2}{Divisibilitat}
% Opcional abans de \input: \def\fmneed{14} (espai mínim abans de cada secció).
% =====================================================================
\makeatletter
\def\fm@logodir{../../logo/pdf}
\providecommand{\fmneed}{10}
\def\fm@cursnom{}

% clau del curs, color i color del text damunt del color
\newcommand{\fm@curs}[3]{\@namedef{fm@color@#1}{#2}\@namedef{fm@text@#1}{#3}}
\fm@curs{1eso}{E6B84A}{navy}
\fm@curs{2eso}{3B82F6}{white}
\fm@curs{3eso}{FF7A50}{white}
\fm@curs{4eso-a}{A78BFA}{white}
\fm@curs{4eso-b}{6D28D9}{white}
\fm@curs{1batx-mi}{9B1C1C}{white}
\fm@curs{1batx-ccss}{EF6B6B}{white}
\fm@curs{2batx-mii}{1E6B47}{white}
\fm@curs{2batx-ccss}{2BB3A3}{white}

\newcommand{\fmcurs}[1]{%
  \@ifundefined{fm@color@#1}{\PackageError{fentmates}{Curs desconegut: #1}{}}{}%
  \edef\fm@tmp{\noexpand\definecolor{curs}{HTML}{\@nameuse{fm@color@#1}}%
               \noexpand\colorlet{curstext}{\@nameuse{fm@text@#1}}}\fm@tmp
  \colorlet{cursfosc}{curs!72!black}%
  \def\fm@logo{\fm@logodir/vertical-#1.pdf}%
  \def\fm@marca{\fm@logodir/marca-#1.pdf}}

% ---------- capçalera de la primera pàgina ----------
% #1 matèria (petit), #2 curs, #3 tipus de document i unitat, #4 títol
\newcommand{\capcalera}[4]{%
  \gdef\fm@cursnom{#2}%
  \par\noindent
  \begin{tcolorbox}[colback=curs,colframe=curs,boxrule=0pt,arc=2.2mm,width=3.5cm,height=2.15cm,
      valign=center,left=3mm,right=2mm,top=1mm,bottom=1mm,nobeforeafter,box align=center]
    \color{curstext}{\footnotesize #1}\par\vspace{1pt}{\fontsize{19}{21}\selectfont\bfseries #2\par}%
  \end{tcolorbox}\hspace{5mm}%
  \begin{minipage}[c]{\dimexpr\linewidth-3.5cm-5mm-2.9cm\relax}
    {\color{cursfosc}\small\bfseries\MakeUppercase{#3}}\par\vspace{2pt}
    {\color{navy}\fontsize{20}{23}\selectfont\bfseries #4\par}
  \end{minipage}\hfill
  \begin{minipage}[c]{2.6cm}\raggedleft\includegraphics[width=2.6cm]{\fm@logo}\end{minipage}\par
  \vspace{2pt}}

% ---------- barres de secció del color del curs ----------
\renewcommand{\seccio}[2]{\par\needspace{\fmneed\baselineskip}\bigskip
  \begin{tcolorbox}[colback=curs,colframe=curs,coltext=curstext,arc=2mm,boxrule=0pt,left=3mm,right=3mm,top=1.2mm,bottom=1.2mm]
  {\large\bfseries #1\quad #2}
  \end{tcolorbox}\vspace{-2pt}}

% ---------- estrelles de dificultat del color del curs (a Batxillerat no se'n posen) ----------
\@ifundefined{estrelles}{}{%
  \renewcommand{\estrelles}[1]{{\color{curs}\ifcase#1\or\faStar\faStar[regular]\faStar[regular]\or\faStar\faStar\faStar[regular]\or\faStar\faStar\faStar\fi}}}

% ---------- capçalera i peu de les pàgines ----------
% primera pàgina: logos institucionals al peu
\fancypagestyle{primera}{\fancyhf{}\renewcommand{\headrulewidth}{0pt}
  \fancyfoot[L]{\raisebox{-0.32\height}{\includegraphics[height=0.78cm]{logo-gva.png}}\hspace{5mm}%
    \raisebox{-0.32\height}{\includegraphics[height=0.78cm]{logo-vermellar.png}}}
  \fancyfoot[R]{\scriptsize\color{muted}\href{https://creativecommons.org/licenses/by-nc-sa/4.0/deed.ca}{\ccbyncsa}\ \ CC BY-NC-SA 4.0 · fentmates.org\qquad\small\thepage}}
% la resta: el símbol a l'esquerra i el curs a la dreta, amb el títol al centre com abans
\fancyhead[L]{\raisebox{-0.3\height}{\includegraphics[height=0.62cm]{\fm@marca}}}
\fancyhead[R]{\small\bfseries\color{cursfosc}\fm@cursnom}
\renewcommand{\headrule}{\color{curs!45}\hrule height \headrulewidth}
\makeatother

%   \fmcurs{1eso}
% i, al començament del document:
%   \capcalera{Matemàtiques}{1r ESO}{Fitxa d'exercicis · Unitat 2}{Divisibilitat}
% Opcional abans de \input: \def\fmneed{14} (espai mínim abans de cada secció).
% =====================================================================
\makeatletter
\def\fm@logodir{../../logo/pdf}
\providecommand{\fmneed}{10}
\def\fm@cursnom{}

% clau del curs, color i color del text damunt del color
\newcommand{\fm@curs}[3]{\@namedef{fm@color@#1}{#2}\@namedef{fm@text@#1}{#3}}
\fm@curs{1eso}{E6B84A}{navy}
\fm@curs{2eso}{3B82F6}{white}
\fm@curs{3eso}{FF7A50}{white}
\fm@curs{4eso-a}{A78BFA}{white}
\fm@curs{4eso-b}{6D28D9}{white}
\fm@curs{1batx-mi}{9B1C1C}{white}
\fm@curs{1batx-ccss}{EF6B6B}{white}
\fm@curs{2batx-mii}{1E6B47}{white}
\fm@curs{2batx-ccss}{2BB3A3}{white}

\newcommand{\fmcurs}[1]{%
  \@ifundefined{fm@color@#1}{\PackageError{fentmates}{Curs desconegut: #1}{}}{}%
  \edef\fm@tmp{\noexpand\definecolor{curs}{HTML}{\@nameuse{fm@color@#1}}%
               \noexpand\colorlet{curstext}{\@nameuse{fm@text@#1}}}\fm@tmp
  \colorlet{cursfosc}{curs!72!black}%
  \def\fm@logo{\fm@logodir/vertical-#1.pdf}%
  \def\fm@marca{\fm@logodir/marca-#1.pdf}}

% ---------- capçalera de la primera pàgina ----------
% #1 matèria (petit), #2 curs, #3 tipus de document i unitat, #4 títol
\newcommand{\capcalera}[4]{%
  \gdef\fm@cursnom{#2}%
  \par\noindent
  \begin{tcolorbox}[colback=curs,colframe=curs,boxrule=0pt,arc=2.2mm,width=3.5cm,height=2.15cm,
      valign=center,left=3mm,right=2mm,top=1mm,bottom=1mm,nobeforeafter,box align=center]
    \color{curstext}{\footnotesize #1}\par\vspace{1pt}{\fontsize{19}{21}\selectfont\bfseries #2\par}%
  \end{tcolorbox}\hspace{5mm}%
  \begin{minipage}[c]{\dimexpr\linewidth-3.5cm-5mm-2.9cm\relax}
    {\color{cursfosc}\small\bfseries\MakeUppercase{#3}}\par\vspace{2pt}
    {\color{navy}\fontsize{20}{23}\selectfont\bfseries #4\par}
  \end{minipage}\hfill
  \begin{minipage}[c]{2.6cm}\raggedleft\includegraphics[width=2.6cm]{\fm@logo}\end{minipage}\par
  \vspace{2pt}}

% ---------- barres de secció del color del curs ----------
\renewcommand{\seccio}[2]{\par\needspace{\fmneed\baselineskip}\bigskip
  \begin{tcolorbox}[colback=curs,colframe=curs,coltext=curstext,arc=2mm,boxrule=0pt,left=3mm,right=3mm,top=1.2mm,bottom=1.2mm]
  {\large\bfseries #1\quad #2}
  \end{tcolorbox}\vspace{-2pt}}

% ---------- estrelles de dificultat del color del curs (a Batxillerat no se'n posen) ----------
\@ifundefined{estrelles}{}{%
  \renewcommand{\estrelles}[1]{{\color{curs}\ifcase#1\or\faStar\faStar[regular]\faStar[regular]\or\faStar\faStar\faStar[regular]\or\faStar\faStar\faStar\fi}}}

% ---------- capçalera i peu de les pàgines ----------
% primera pàgina: logos institucionals al peu
\fancypagestyle{primera}{\fancyhf{}\renewcommand{\headrulewidth}{0pt}
  \fancyfoot[L]{\raisebox{-0.32\height}{\includegraphics[height=0.78cm]{logo-gva.png}}\hspace{5mm}%
    \raisebox{-0.32\height}{\includegraphics[height=0.78cm]{logo-vermellar.png}}}
  \fancyfoot[R]{\scriptsize\color{muted}\href{https://creativecommons.org/licenses/by-nc-sa/4.0/deed.ca}{\ccbyncsa}\ \ CC BY-NC-SA 4.0 · fentmates.org\qquad\small\thepage}}
% la resta: el símbol a l'esquerra i el curs a la dreta, amb el títol al centre com abans
\fancyhead[L]{\raisebox{-0.3\height}{\includegraphics[height=0.62cm]{\fm@marca}}}
\fancyhead[R]{\small\bfseries\color{cursfosc}\fm@cursnom}
\renewcommand{\headrule}{\color{curs!45}\hrule height \headrulewidth}
\makeatother

%   \fmcurs{1eso}
% i, al començament del document:
%   \capcalera{Matemàtiques}{1r ESO}{Fitxa d'exercicis · Unitat 2}{Divisibilitat}
% Opcional abans de \input: \def\fmneed{14} (espai mínim abans de cada secció).
% =====================================================================
\makeatletter
\def\fm@logodir{../../logo/pdf}
\providecommand{\fmneed}{10}
\def\fm@cursnom{}

% clau del curs, color i color del text damunt del color
\newcommand{\fm@curs}[3]{\@namedef{fm@color@#1}{#2}\@namedef{fm@text@#1}{#3}}
\fm@curs{1eso}{E6B84A}{navy}
\fm@curs{2eso}{3B82F6}{white}
\fm@curs{3eso}{FF7A50}{white}
\fm@curs{4eso-a}{A78BFA}{white}
\fm@curs{4eso-b}{6D28D9}{white}
\fm@curs{1batx-mi}{9B1C1C}{white}
\fm@curs{1batx-ccss}{EF6B6B}{white}
\fm@curs{2batx-mii}{1E6B47}{white}
\fm@curs{2batx-ccss}{2BB3A3}{white}

\newcommand{\fmcurs}[1]{%
  \@ifundefined{fm@color@#1}{\PackageError{fentmates}{Curs desconegut: #1}{}}{}%
  \edef\fm@tmp{\noexpand\definecolor{curs}{HTML}{\@nameuse{fm@color@#1}}%
               \noexpand\colorlet{curstext}{\@nameuse{fm@text@#1}}}\fm@tmp
  \colorlet{cursfosc}{curs!72!black}%
  \def\fm@logo{\fm@logodir/vertical-#1.pdf}%
  \def\fm@marca{\fm@logodir/marca-#1.pdf}}

% ---------- capçalera de la primera pàgina ----------
% #1 matèria (petit), #2 curs, #3 tipus de document i unitat, #4 títol
\newcommand{\capcalera}[4]{%
  \gdef\fm@cursnom{#2}%
  \par\noindent
  \begin{tcolorbox}[colback=curs,colframe=curs,boxrule=0pt,arc=2.2mm,width=3.5cm,height=2.15cm,
      valign=center,left=3mm,right=2mm,top=1mm,bottom=1mm,nobeforeafter,box align=center]
    \color{curstext}{\footnotesize #1}\par\vspace{1pt}{\fontsize{19}{21}\selectfont\bfseries #2\par}%
  \end{tcolorbox}\hspace{5mm}%
  \begin{minipage}[c]{\dimexpr\linewidth-3.5cm-5mm-2.9cm\relax}
    {\color{cursfosc}\small\bfseries\MakeUppercase{#3}}\par\vspace{2pt}
    {\color{navy}\fontsize{20}{23}\selectfont\bfseries #4\par}
  \end{minipage}\hfill
  \begin{minipage}[c]{2.6cm}\raggedleft\includegraphics[width=2.6cm]{\fm@logo}\end{minipage}\par
  \vspace{2pt}}

% ---------- barres de secció del color del curs ----------
\renewcommand{\seccio}[2]{\par\needspace{\fmneed\baselineskip}\bigskip
  \begin{tcolorbox}[colback=curs,colframe=curs,coltext=curstext,arc=2mm,boxrule=0pt,left=3mm,right=3mm,top=1.2mm,bottom=1.2mm]
  {\large\bfseries #1\quad #2}
  \end{tcolorbox}\vspace{-2pt}}

% ---------- estrelles de dificultat del color del curs (a Batxillerat no se'n posen) ----------
\@ifundefined{estrelles}{}{%
  \renewcommand{\estrelles}[1]{{\color{curs}\ifcase#1\or\faStar\faStar[regular]\faStar[regular]\or\faStar\faStar\faStar[regular]\or\faStar\faStar\faStar\fi}}}

% ---------- capçalera i peu de les pàgines ----------
% primera pàgina: logos institucionals al peu
\fancypagestyle{primera}{\fancyhf{}\renewcommand{\headrulewidth}{0pt}
  \fancyfoot[L]{\raisebox{-0.32\height}{\includegraphics[height=0.78cm]{logo-gva.png}}\hspace{5mm}%
    \raisebox{-0.32\height}{\includegraphics[height=0.78cm]{logo-vermellar.png}}}
  \fancyfoot[R]{\scriptsize\color{muted}\href{https://creativecommons.org/licenses/by-nc-sa/4.0/deed.ca}{\ccbyncsa}\ \ CC BY-NC-SA 4.0 · fentmates.org\qquad\small\thepage}}
% la resta: el símbol a l'esquerra i el curs a la dreta, amb el títol al centre com abans
\fancyhead[L]{\raisebox{-0.3\height}{\includegraphics[height=0.62cm]{\fm@marca}}}
\fancyhead[R]{\small\bfseries\color{cursfosc}\fm@cursnom}
\renewcommand{\headrule}{\color{curs!45}\hrule height \headrulewidth}
\makeatother

\fmcurs{2eso}

\begin{document}
\thispagestyle{primera}

\capcalera{Matemáticas}{2\textsuperscript{o} ESO}{Ficha de ejercicios · Unidad 9}{Cuerpos geométricos.\\ Áreas y volúmenes}
\vspace{10pt}
Nombre y apellidos: \rule{6.6cm}{0.4pt}\hfill Grupo: \rule{1.6cm}{0.4pt}\hfill Fecha: \rule{2.2cm}{0.4pt}
\vspace{8pt}

\begin{tcolorbox}[colback=soft,colframe=soft,arc=2mm,left=3mm,right=3mm,top=1.5mm,bottom=1.5mm]
\small Ejercicios graduados por niveles de dificultad:\quad \estrelles{1}\ básico \qquad \estrelles{2}\ medio \qquad \estrelles{3}\ avanzado.\\[2pt] Resuélvelos en el cuaderno mostrando todos los pasos y haz un dibujo de cada cuerpo. Si un resultado no es exacto, redondéalo a las centésimas y usa la tecla $\pi$ de la calculadora.
\end{tcolorbox}

% =====================================================
\seccio{1}{Poliedros. Fórmula de Euler. Poliedros regulares}

\ex{1} Copia y completa la tabla. Comprueba en cada caso la fórmula de Euler.\par\smallskip\begin{center}\renewcommand{\arraystretch}{1.25}\begin{tabular}{|l|C{1.6cm}|C{1.6cm}|C{1.6cm}|}\hline Poliedro & Caras & Vértices & Aristas \\\hline prisma pentagonal & & &  \\\hline pirámide hexagonal & & &  \\\hline prisma triangular & & &  \\\hline pirámide cuadrangular & & &  \\\hline ortoedro & & &  \\\hline\end{tabular}\end{center}\recorda{En un poliedro convexo, \textbf{caras} + \textbf{vértices} = \textbf{aristas} + 2: \ $C+V=A+2$ (fórmula de Euler). Un prisma de base de $n$ lados tiene $n+2$ caras, $2n$ vértices y $3n$ aristas; una pirámide, $n+1$ caras, $n+1$ vértices y $2n$ aristas.}

\ex{1} Usa la fórmula de Euler para calcular lo que falta en cada poliedro convexo:\begin{tasks}(3)
  \task $C=6$, \ $V=8$, \ $A=\ ?$
  \task $V=12$, \ $A=18$, \ $C=\ ?$
  \task $C=7$, \ $A=12$, \ $V=\ ?$
\end{tasks}

\ex{1} ¿Cuáles de estos poliedros son regulares? Justifícalo.\begin{tasks}(3)
  \task tetraedro
  \task prisma triangular
  \task cubo
  \task pirámide cuadrangular
  \task dodecaedro
  \task ortoedro
\end{tasks}

\ex{2} Escribe, para cada uno de los cinco poliedros regulares, cuántas caras tiene, qué forma tienen y cuántas caras se encuentran en cada vértice.

\ex{2} Un prisma tiene 18 aristas. ¿Qué polígono tiene por base? ¿Cuántas caras y cuántos vértices tiene?

\ex{2} Una pirámide tiene 9 caras. ¿Qué polígono tiene por base? ¿Cuántos vértices y aristas tiene?

\ex{3} ¿Puede existir un poliedro convexo con: \ a) 8 caras, 12 vértices y 18 aristas? \ b) 6 caras, 7 vértices y 12 aristas?

\ex{3} Explica por qué no puede haber ningún poliedro regular con caras hexagonales.


% =====================================================
\seccio{2}{Prismas: elementos y área}

\ex{1} \begin{minipage}[t]{0.50\linewidth}Calcula el área total:\begin{tasks}(1)
  \task de un cubo de 5 cm de arista;
  \task del ortoedro de la figura.
\end{tasks}\end{minipage}\hfill\begin{minipage}[t]{0.36\linewidth}\vspace{0pt}\centering \begin{tikzpicture}[x=0.42cm,y=0.42cm,baseline=(current bounding box.center),line join=round]
\path[fill=teal!14,draw=none] (0.000,2.782) -- (0.000,0.000) -- (6.181,-1.231) -- (6.181,1.550) -- cycle;
\path[fill=teal!14,draw=none] (6.181,1.550) -- (6.181,-1.231) -- (8.059,0.092) -- (8.059,2.874) -- cycle;
\path[fill=teal!14,draw=none] (1.878,4.105) -- (0.000,2.782) -- (6.181,1.550) -- (8.059,2.874) -- cycle;
\draw[dashed,draw=muted] (0.000,0.000) -- (1.878,1.323);
\draw[dashed,draw=muted] (1.878,1.323) -- (8.059,0.092);
\draw[dashed,draw=muted] (1.878,1.323) -- (1.878,4.105);
\draw[thick] (0.000,0.000) -- (6.181,-1.231);
\draw[thick] (6.181,-1.231) -- (8.059,0.092);
\draw[thick] (0.000,0.000) -- (0.000,2.782);
\draw[thick] (6.181,-1.231) -- (6.181,1.550);
\draw[thick] (0.000,2.782) -- (6.181,1.550);
\draw[thick] (0.000,2.782) -- (1.878,4.105);
\draw[thick] (8.059,0.092) -- (8.059,2.874);
\draw[thick] (6.181,1.550) -- (8.059,2.874);
\draw[thick] (1.878,4.105) -- (8.059,2.874);
\node[font=\small,inner sep=2pt,below] at (3.090,-0.616) {$7$\,cm};
\node[font=\small,inner sep=2pt,right] at (7.120,-0.570) {$4$\,cm};
\node[font=\small,inner sep=2pt,right] at (6.181,0.160) {$3$\,cm};
\end{tikzpicture}\end{minipage}\par\recorda{Área de un prisma recto: $A_L=P_{\text{base}}\cdot h$, \ $A_T=A_L+2\,A_{\text{base}}$. Ortoedro de aristas $a,\,b,\,c$: $A_T=2(ab+ac+bc)$ y la diagonal es $D=\sqrt{a^2+b^2+c^2}$.}

\ex{1} Calcula la diagonal de un ortoedro de 6 cm × 2 cm × 3 cm.

\ex{1} Un prisma cuadrangular regular tiene 4 cm de arista de la base y 10 cm de altura. Calcula su área lateral y su área total.

\ex{2} La base de un prisma recto es un triángulo rectángulo de catetos 5 cm y 12 cm, y la altura del prisma es de 10 cm. Calcula su área total.

\ex{2} Calcula el área total de un prisma hexagonal regular de 8 cm de arista de la base y 5 cm de altura.

\ex{2} Queremos hacer una caja sin tapa de 30 cm × 20 cm de base y 15 cm de altura. ¿Cuánto cartón necesitamos?

\ex{3} Un cubo tiene un área total de 216 cm². ¿Cuánto mide la arista? ¿Y la diagonal?

\ex{3} Queremos pintar las cuatro paredes de una habitación de 4 m × 3,5 m y 2,5 m de altura. La puerta mide 0,8 m × 2 m y la ventana 1,2 m × 1 m. Si pintar cuesta 4 € el metro cuadrado, ¿cuánto nos costará?


% =====================================================
\seccio{3}{Pirámides: elementos y área}

\ex{1} \begin{minipage}[t]{0.54\linewidth}La pirámide cuadrangular regular de la figura tiene 16 cm de arista de la base y 6 cm de altura. Calcula su apotema y su área total.\end{minipage}\hfill\begin{minipage}[t]{0.32\linewidth}\vspace{0pt}\centering \begin{tikzpicture}[x=0.6cm,y=0.6cm,baseline=(current bounding box.center),line join=round]
\path[fill=teal!14,draw=none] (0.000,2.967) -- (-2.434,-0.279) -- (0.744,-0.912) -- cycle;
\path[fill=teal!14,draw=none] (0.000,2.967) -- (0.744,-0.912) -- (2.434,0.279) -- cycle;
\draw[dashed,draw=muted] (-2.434,-0.279) -- (-0.744,0.912);
\draw[dashed,draw=muted] (2.434,0.279) -- (-0.744,0.912);
\draw[dashed,draw=muted] (-0.744,0.912) -- (0.000,2.967);
\draw[thick] (-2.434,-0.279) -- (0.744,-0.912);
\draw[thick] (0.744,-0.912) -- (2.434,0.279);
\draw[thick] (-2.434,-0.279) -- (0.000,2.967);
\draw[thick] (0.744,-0.912) -- (0.000,2.967);
\draw[thick] (2.434,0.279) -- (0.000,2.967);
\node[font=\small,inner sep=2pt,below] at (-0.845,-0.595) {$16$\,cm};
\draw[dashed,draw=accent,thick] (0.000,0.000) -- (0.000,2.967);
\draw[dashed,draw=accent,thick] (0.000,0.000) -- (1.589,-0.317);
\draw[draw=navy,thick] (1.589,-0.317) -- (0.000,2.967);
\fill (0.000,0.000) circle (1.3pt);
\node[font=\small,inner sep=2pt,left,text=accent] at (0.000,1.483) {$6$\,cm};
\end{tikzpicture}\end{minipage}\par\recorda{En una pirámide regular, la altura $h$, la apotema de la base $a_b$ y la apotema de la pirámide $a_p$ forman un triángulo rectángulo: $a_p^2=h^2+a_b^2$. \ $A_L=\frac{P_{\text{base}}\cdot a_p}{2}$, \ $A_T=A_L+A_{\text{base}}$.}

\ex{1} Una pirámide hexagonal regular tiene 6 cm de arista de la base y 10 cm de apotema. Calcula su área lateral.

\ex{1} Una pirámide triangular regular tiene 8 cm de arista de la base y 12 cm de apotema. Calcula su área lateral.

\ex{2} Una pirámide cuadrangular regular tiene 10 cm de arista de la base y 13 cm de apotema. Calcula su altura y su área total.

\ex{2} Una pirámide cuadrangular regular tiene 6 cm de arista de la base y 4 cm de altura. Calcula la apotema, la arista lateral y el área total.

\ex{2} Calcula el área total de un tetraedro regular de 4 cm de arista.

\ex{3} Una tienda de campaña tiene forma de pirámide cuadrangular regular de 2 m de lado de la base y 1,5 m de altura. ¿Cuánta lona hace falta para hacerla, si no tiene suelo?

\ex{3} Calcula el área total de una pirámide hexagonal regular de 4 cm de arista de la base y 6 cm de altura.


% =====================================================
\seccio{4}{Cuerpos de revolución: cilindro y cono}

\ex{1} ¿Qué cuerpo de revolución se obtiene al hacer girar:\begin{tasks}(1)
  \task un rectángulo alrededor de un lado?
  \task un triángulo rectángulo alrededor de un cateto?
  \task un semicírculo alrededor del diámetro?
  \task un triángulo rectángulo alrededor de la hipotenusa?
\end{tasks}

\ex{1} \begin{minipage}[t]{0.60\linewidth}Calcula el área lateral y el área total del cilindro de la figura, de 3 cm de radio y 8 cm de altura.\end{minipage}\hfill\begin{minipage}[t]{0.26\linewidth}\vspace{0pt}\centering \begin{tikzpicture}[x=0.45cm,y=0.45cm,baseline=(current bounding box.center),line join=round]
\path[fill=teal!14,draw=none] (-1.500,3.709) -- (-1.500,0.000) -- (-1.497,-0.037) -- (-1.487,-0.073) -- (-1.471,-0.110) -- (-1.449,-0.145) -- (-1.420,-0.181) -- (-1.386,-0.215) -- (-1.345,-0.249) -- (-1.299,-0.281) -- (-1.247,-0.312) -- (-1.190,-0.342) -- (-1.128,-0.370) -- (-1.061,-0.397) -- (-0.989,-0.422) -- (-0.913,-0.446) -- (-0.833,-0.467) -- (-0.750,-0.487) -- (-0.663,-0.504) -- (-0.574,-0.519) -- (-0.482,-0.532) -- (-0.388,-0.543) -- (-0.293,-0.551) -- (-0.196,-0.557) -- (-0.098,-0.561) -- (-0.000,-0.562) -- (0.098,-0.561) -- (0.196,-0.557) -- (0.293,-0.551) -- (0.388,-0.543) -- (0.482,-0.532) -- (0.574,-0.519) -- (0.663,-0.504) -- (0.750,-0.487) -- (0.833,-0.467) -- (0.913,-0.446) -- (0.989,-0.422) -- (1.061,-0.397) -- (1.128,-0.370) -- (1.190,-0.342) -- (1.247,-0.312) -- (1.299,-0.281) -- (1.345,-0.249) -- (1.386,-0.215) -- (1.420,-0.181) -- (1.449,-0.145) -- (1.471,-0.110) -- (1.487,-0.073) -- (1.497,-0.037) -- (1.500,-0.000) -- (1.500,3.709) -- (1.500,3.709) -- (1.497,3.672) -- (1.487,3.635) -- (1.471,3.599) -- (1.449,3.563) -- (1.420,3.528) -- (1.386,3.494) -- (1.345,3.460) -- (1.299,3.428) -- (1.247,3.397) -- (1.190,3.367) -- (1.128,3.338) -- (1.061,3.311) -- (0.989,3.286) -- (0.913,3.263) -- (0.833,3.242) -- (0.750,3.222) -- (0.663,3.205) -- (0.574,3.190) -- (0.482,3.177) -- (0.388,3.166) -- (0.293,3.158) -- (0.196,3.152) -- (0.098,3.148) -- (-0.000,3.147) -- (-0.098,3.148) -- (-0.196,3.152) -- (-0.293,3.158) -- (-0.388,3.166) -- (-0.482,3.177) -- (-0.574,3.190) -- (-0.663,3.205) -- (-0.750,3.222) -- (-0.833,3.242) -- (-0.913,3.263) -- (-0.989,3.286) -- (-1.061,3.311) -- (-1.128,3.338) -- (-1.190,3.367) -- (-1.247,3.397) -- (-1.299,3.428) -- (-1.345,3.460) -- (-1.386,3.494) -- (-1.420,3.528) -- (-1.449,3.563) -- (-1.471,3.599) -- (-1.487,3.635) -- (-1.497,3.672) -- (-1.500,3.709) -- cycle;
\draw[dashed,draw=muted] (1.500,0.000) -- (1.497,0.037) -- (1.487,0.073) -- (1.471,0.110) -- (1.449,0.145) -- (1.420,0.181) -- (1.386,0.215) -- (1.345,0.249) -- (1.299,0.281) -- (1.247,0.312) -- (1.190,0.342) -- (1.128,0.370) -- (1.061,0.397) -- (0.989,0.422) -- (0.913,0.446) -- (0.833,0.467) -- (0.750,0.487) -- (0.663,0.504) -- (0.574,0.519) -- (0.482,0.532) -- (0.388,0.543) -- (0.293,0.551) -- (0.196,0.557) -- (0.098,0.561) -- (0.000,0.562) -- (-0.098,0.561) -- (-0.196,0.557) -- (-0.293,0.551) -- (-0.388,0.543) -- (-0.482,0.532) -- (-0.574,0.519) -- (-0.663,0.504) -- (-0.750,0.487) -- (-0.833,0.467) -- (-0.913,0.446) -- (-0.989,0.422) -- (-1.061,0.397) -- (-1.128,0.370) -- (-1.190,0.342) -- (-1.247,0.312) -- (-1.299,0.281) -- (-1.345,0.249) -- (-1.386,0.215) -- (-1.420,0.181) -- (-1.449,0.145) -- (-1.471,0.110) -- (-1.487,0.073) -- (-1.497,0.037) -- (-1.500,0.000);
\draw[thick] (-1.500,0.000) -- (-1.497,-0.037) -- (-1.487,-0.073) -- (-1.471,-0.110) -- (-1.449,-0.145) -- (-1.420,-0.181) -- (-1.386,-0.215) -- (-1.345,-0.249) -- (-1.299,-0.281) -- (-1.247,-0.312) -- (-1.190,-0.342) -- (-1.128,-0.370) -- (-1.061,-0.397) -- (-0.989,-0.422) -- (-0.913,-0.446) -- (-0.833,-0.467) -- (-0.750,-0.487) -- (-0.663,-0.504) -- (-0.574,-0.519) -- (-0.482,-0.532) -- (-0.388,-0.543) -- (-0.293,-0.551) -- (-0.196,-0.557) -- (-0.098,-0.561) -- (-0.000,-0.562) -- (0.098,-0.561) -- (0.196,-0.557) -- (0.293,-0.551) -- (0.388,-0.543) -- (0.482,-0.532) -- (0.574,-0.519) -- (0.663,-0.504) -- (0.750,-0.487) -- (0.833,-0.467) -- (0.913,-0.446) -- (0.989,-0.422) -- (1.061,-0.397) -- (1.128,-0.370) -- (1.190,-0.342) -- (1.247,-0.312) -- (1.299,-0.281) -- (1.345,-0.249) -- (1.386,-0.215) -- (1.420,-0.181) -- (1.449,-0.145) -- (1.471,-0.110) -- (1.487,-0.073) -- (1.497,-0.037) -- (1.500,-0.000);
\draw[thick] (-1.500,0.000) -- (-1.500,3.709);
\draw[thick] (1.500,0.000) -- (1.500,3.709);
\draw[fill=teal!26,draw=black,thick] (1.500,3.709) -- (1.487,3.782) -- (1.449,3.854) -- (1.386,3.924) -- (1.299,3.990) -- (1.190,4.051) -- (1.061,4.106) -- (0.913,4.155) -- (0.750,4.195) -- (0.574,4.228) -- (0.388,4.251) -- (0.196,4.266) -- (0.000,4.271) -- (-0.196,4.266) -- (-0.388,4.251) -- (-0.574,4.228) -- (-0.750,4.195) -- (-0.913,4.155) -- (-1.061,4.106) -- (-1.190,4.051) -- (-1.299,3.990) -- (-1.386,3.924) -- (-1.449,3.854) -- (-1.487,3.782) -- (-1.500,3.709) -- (-1.487,3.635) -- (-1.449,3.563) -- (-1.386,3.494) -- (-1.299,3.428) -- (-1.190,3.367) -- (-1.061,3.311) -- (-0.913,3.263) -- (-0.750,3.222) -- (-0.574,3.190) -- (-0.388,3.166) -- (-0.196,3.152) -- (-0.000,3.147) -- (0.196,3.152) -- (0.388,3.166) -- (0.574,3.190) -- (0.750,3.222) -- (0.913,3.263) -- (1.061,3.311) -- (1.190,3.367) -- (1.299,3.428) -- (1.386,3.494) -- (1.449,3.563) -- (1.487,3.635) -- (1.500,3.709) -- cycle;
\draw[dashed,draw=accent,thick] (0.000,0.000) -- (1.500,0.000);
\fill (0.000,0.000) circle (1.3pt);
\node[font=\small,inner sep=2pt,above,text=accent] at (0.750,0.000) {$3$\,cm};
\node[font=\small,inner sep=2pt,right] at (1.500,1.854) {$8$\,cm};
\end{tikzpicture}\end{minipage}\par\recorda{Cilindro: $A_L=2\pi r h$, \ $A_T=2\pi r h+2\pi r^2$. \ Cono: $A_L=\pi r g$, \ $A_T=\pi r g+\pi r^2$, donde la generatriz cumple $g^2=h^2+r^2$.}

\ex{1} Un cono tiene 6 cm de radio y 10 cm de generatriz. Calcula su área lateral y su área total.

\ex{2} Un cono tiene 16 cm de diámetro y 15 cm de altura. Calcula la generatriz y el área total.

\ex{2} La etiqueta de una lata cilíndrica de 8 cm de diámetro y 11 cm de altura cubre toda la superficie lateral. ¿Qué área tiene?

\ex{2} Un cono tiene 25 cm de generatriz y 7 cm de radio. ¿Qué altura tiene?

\ex{3} Un cilindro de 10 cm de altura tiene un área lateral de 314,16 cm². ¿Qué radio tiene? Calcula su área total.

\ex{3} Queremos hacer vasos cilíndricos de cartulina, sin tapa, de 3 cm de radio y 9 cm de altura. ¿Cuánta cartulina hace falta para cada vaso? ¿Cuántos vasos, como máximo, podríamos hacer con 1 m² de cartulina (sin contar lo que se pierde al recortar)?


% =====================================================
\seccio{5}{La esfera y su área}

\ex{1} Calcula el área de:\begin{tasks}(2)
  \task una esfera de 7 cm de radio;
  \task una esfera de 10 cm de diámetro.
\end{tasks}\recorda{El área de la esfera es la de cuatro círculos máximos: $A=4\pi r^2$. Cuidado: si te dan el diámetro, el radio es la mitad.}

\ex{1} Una esfera tiene un área de $36\pi$ cm². ¿Qué radio tiene?

\ex{2} Una pelota de tenis mide 6,5 cm de diámetro. ¿Cuánta superficie tiene?

\ex{2} Queremos pintar por fuera una cúpula semiesférica de 6 m de radio. ¿Qué superficie tenemos que pintar?

\ex{2} La circunferencia máxima de un balón mide 50 cm. Calcula su radio y su área.

\ex{3} Una esfera está metida dentro de un cubo de 10 cm de arista y toca todas las caras. ¿Cuál es el área de la esfera? ¿Qué porcentaje representa del área del cubo?

\ex{3} Una esfera de 4 cm de radio está metida en un cilindro que la envuelve exactamente (radio 4 cm y altura 8 cm). Comprueba que el área de la esfera es igual al área lateral del cilindro. ¿Qué fracción del área total del cilindro representa?


% =====================================================
\seccio{6}{Volumen. Prismas y cilindros}

\ex{1} Expresa en la unidad que se indica:\begin{tasks}(3)
  \task 5 m³ \ en dm³
  \task 2,5 dm³ \ en cm³
  \task 7500 cm³ \ en dm³
  \task 0,35 m³ \ en litros
  \task 3 L \ en cm³
  \task 250 mL \ en dm³
\end{tasks}\recorda{Cada unidad de volumen es 1000 veces la siguiente: 1 m³ = 1000 dm³, 1 dm³ = 1000 cm³. \ Y la clave de la capacidad: 1 dm³ = 1 L, \ 1 cm³ = 1 mL, \ 1 m³ = 1000 L.}

\ex{1} Calcula el volumen de un cubo de 6 cm de arista y el de un ortoedro de 10 cm × 4 cm × 5 cm.

\ex{1} Calcula el volumen de un cilindro de 2 cm de radio y 10 cm de altura.

\ex{2} Un prisma tiene de base un triángulo rectángulo de catetos 3 cm y 4 cm, y 20 cm de altura. ¿Qué volumen tiene?

\ex{2} Un acuario mide 60 cm de largo, 30 cm de ancho y 40 cm de alto. ¿Cuántos litros caben? ¿Cuántos ponemos si lo llenamos hasta 35 cm de altura?

\ex{2} Un depósito cilíndrico tiene 1 m de radio y 3 m de altura. ¿Cuántos litros de agua caben?

\ex{3} Calcula el volumen de un prisma hexagonal regular de 2 cm de arista de la base y 10 cm de altura.

\ex{3} Un tetrabrik de 1 L de leche tiene una base rectangular de 6 cm × 9,5 cm. ¿Qué altura debe tener como mínimo?


% =====================================================
\seccio{7}{Volumen de pirámides, conos y esferas}

\ex{1} Calcula el volumen de una pirámide cuadrangular de 9 cm de arista de la base y 10 cm de altura.\recorda{$V_{\text{pirámide}}=\frac{A_{\text{base}}\cdot h}{3}$, \ $V_{\text{cono}}=\frac{\pi r^2h}{3}$ (un tercio del prisma o del cilindro de la misma base y altura), \ $V_{\text{esfera}}=\frac43\pi r^3$.}

\ex{1} Calcula el volumen de un cono de 4 cm de radio y 9 cm de altura.

\ex{1} Calcula el volumen de una esfera de 3 cm de radio y el de una de 10 cm de radio.

\ex{2} Un cono tiene 15 cm de generatriz y 9 cm de radio. Calcula su altura y su volumen.

\ex{2} Una pirámide cuadrangular regular tiene 8 cm de arista de la base y 5 cm de apotema. Calcula su altura y su volumen.

\ex{2} Un balón de baloncesto mide 24 cm de diámetro. ¿Cuántos litros de aire contiene?

\ex{3} Una esfera tiene un volumen de $288\pi$ cm³. ¿Qué radio tiene? Calcula su área.

\ex{3} Un cilindro de 3 cm de radio y 6 cm de altura está lleno de agua. Metemos una esfera de 3 cm de radio (que cabe justa). ¿Qué fracción del agua se derrama?


% =====================================================
\seccio{8}{Cuerpos compuestos y problemas}

\ex{1} \begin{minipage}[t]{0.66\linewidth}La pieza de la figura es un cubo de 5 cm de arista con una pirámide de 6 cm de altura encima. Calcula su volumen.\end{minipage}\hfill\begin{minipage}[t]{0.2\linewidth}\vspace{0pt}\centering \begin{tikzpicture}[x=0.3cm,y=0.3cm,baseline=(current bounding box.center),line join=round]
\path[fill=teal!14,draw=none] (0.000,4.636) -- (0.000,0.000) -- (4.415,-0.879) -- (4.415,3.757) -- cycle;
\path[fill=teal!14,draw=none] (4.415,3.757) -- (4.415,-0.879) -- (6.762,0.774) -- (6.762,5.410) -- cycle;
\path[fill=teal!14,draw=none] (3.381,10.586) -- (0.000,4.636) -- (4.415,3.757) -- cycle;
\path[fill=teal!14,draw=none] (3.381,10.586) -- (4.415,3.757) -- (6.762,5.410) -- cycle;
\draw[dashed,draw=muted] (0.000,0.000) -- (2.347,1.654);
\draw[dashed,draw=muted] (2.347,1.654) -- (6.762,0.774);
\draw[dashed,draw=muted] (2.347,1.654) -- (2.347,6.290);
\draw[dashed,draw=muted] (0.000,4.636) -- (2.347,6.290);
\draw[dashed,draw=muted] (2.347,6.290) -- (6.762,5.410);
\draw[dashed,draw=muted] (2.347,6.290) -- (3.381,10.586);
\draw[thick] (0.000,0.000) -- (4.415,-0.879);
\draw[thick] (4.415,-0.879) -- (6.762,0.774);
\draw[thick] (0.000,0.000) -- (0.000,4.636);
\draw[thick] (4.415,-0.879) -- (4.415,3.757);
\draw[thick] (0.000,4.636) -- (4.415,3.757);
\draw[thick] (6.762,0.774) -- (6.762,5.410);
\draw[thick] (4.415,3.757) -- (6.762,5.410);
\draw[thick] (0.000,4.636) -- (3.381,10.586);
\draw[thick] (4.415,3.757) -- (3.381,10.586);
\draw[thick] (6.762,5.410) -- (3.381,10.586);
\end{tikzpicture}\end{minipage}\par\recorda{Para calcular el volumen de un cuerpo compuesto, lo descomponemos en cuerpos simples y sumamos los volúmenes (o los restamos, si hay agujeros). Para el área, solo sumamos las superficies que quedan por fuera.}

\ex{2} Un helado es un cono de 3 cm de radio y 12 cm de altura, lleno de helado, con una semiesfera de helado encima. ¿Cuánto helado hay?

\ex{2} Un silo es un cilindro de 3 m de radio y 8 m de altura con un tejado cónico de 4 m de altura. ¿Qué volumen tiene?

\ex{2} Un tubo de 2 m de largo tiene 6 cm de radio exterior y 5 cm de radio interior. ¿Qué volumen de material tiene?

\ex{2} Un cubo de 12 cm de arista tiene un agujero cilíndrico de 2 cm de radio que lo atraviesa de lado a lado. ¿Qué volumen tiene?

\ex{3} Una piscina mide 25 m de largo, 10 m de ancho y 2 m de profundidad. a) ¿Cuántos metros cuadrados de baldosas hacen falta para cubrir el fondo y las paredes? b) ¿Cuántos litros de agua caben? c) ¿Cuánto tardará en llenarse con una bomba que echa 400 L por minuto?

\ex{3} Tres pelotas de tenis de 3,25 cm de radio se venden en un tubo cilíndrico en el que caben justas (radio 3,25 cm y altura 19,5 cm). ¿Qué fracción del tubo ocupan las pelotas? ¿Qué volumen queda vacío?

\ex{3} Un vaso cilíndrico de 4 cm de radio tiene agua hasta 10 cm de altura. Echamos 5 bolas de 1 cm de radio, que quedan cubiertas. ¿Hasta qué altura llega ahora el agua?

\vfill
\begin{center}\color{muted}\small\faCheckCircle[regular]\ \ ¡Ánimo! Si te cuesta algún ejercicio, repasa los apuntes interactivos de la unidad.\end{center}

\end{document}
